\phantomsection\label{fig:vol01:western-zhou:timeline}
\section{时间线：从牧野到东迁}
\begin{center}
\begin{tikzpicture}[x=.021cm,y=1cm]
  \draw[line width=.8pt,color=timeline] (0,0) -- (275,0);
  \foreach \x in {0,6,69,205,275} {
    \draw[color=dynasty,line width=.8pt] (\x,-.11) -- (\x,.11);
  }
  \node[above,align=center,font=\small\color{ink}] at (0,0) {约前1046\\牧野};
  \node[below,align=center,font=\small\color{ink}] at (6,0) {约前1040\\东征};
  \node[above,align=center,font=\small\color{ink}] at (69,0) {前977\\昭王南征};
  \node[below,align=center,font=\small\color{ink}] at (205,0) {前841\\共和};
  \node[above,align=center,font=\small\color{ink}] at (275,0) {前771\\镐京陷落};
\end{tikzpicture}

\smallskip
\small\textit{图\quad 西周关键节点时间线示意：以《史记·周本纪》《竹书纪年》、金文断代和现代断代研究的对照为线索，标出牧野、东征、昭王南征、共和与镐京陷落。横轴仅表示所列年代的相对间距，不表示事件持续时间、因果强度或逐年材料连续性；所列日期是对照用年代锚点，仍有断代与换算异说。传世文献距事件较远，金文与考古亦不能补足逐年记录，故本图不裁定共和执政者、各次征伐的日期或路线，也不将未安全换算的王年补作连续年表。}
\end{center}
